% 两个三角形{1,2,3}、{4,5,6}由边{3,4}相连：割点3,4，桥{3,4}；链1-2-3与三角形对照。
\begin{tikzpicture}[panelmark/.style={font=\figurefont\fontsize{8}{10}\selectfont,text=FigureInk,anchor=north,inner sep=0pt,outer sep=0pt},
>=Stealth,
  v/.style={circle,draw=FigureStroke,fill=FigureBlueFill,line width=.8pt,minimum size=6.5mm,inner sep=0pt},
  vc/.style={v,draw=FigureRed,fill=FigureRedFill},
  e/.style={draw=FigureStroke,line width=.8pt},
  br/.style={draw=FigureRed,line width=1.6pt},
  note/.style={font=\FigureNoteFont\fontsize{7.5}{9}\selectfont,text=FigureMuted,align=center}]
\begin{scope}[local bounding box=panelrow]
\begin{scope}[local bounding box=panela]
  \node[v] (1) at (-2.1,.7) {$1$}; \node[v] (2) at (-2.1,-.7) {$2$}; \node[vc] (3) at (-.9,0) {$3$};
  \node[vc] (4) at (.5,0) {$4$}; \node[v] (5) at (1.7,.7) {$5$}; \node[v] (6) at (1.7,-.7) {$6$};
  \draw[e] (1)--(2)--(3)--(1) (4)--(5)--(6)--(4);
  \draw[br] (3)--(4);

\end{scope}
\begin{scope}[xshift=3.9cm,local bounding box=panelb]
  \node[v] (1) at (0,0) {$1$}; \node[vc] (2) at (1.4,0) {$2$}; \node[v] (3) at (2.8,0) {$3$};
  \draw[br] (1)--(2)--(3);

\end{scope}
\begin{scope}[xshift=9.6cm,local bounding box=panelc]
  \node[v] (1) at (0,.7) {$1$}; \node[v] (2) at (-.75,-.5) {$2$}; \node[v] (3) at (.75,-.5) {$3$};
  \draw[e] (1)--(2)--(3)--(1);

\end{scope}
\end{scope}
% Subfigure marks: 8 pt, centred below each panel; common top 3 mm below this row.
\coordinate (panelrowbase) at ($(panelrow.south)+(0,-3mm)$);
\node[panelmark] at (panela.center |- panelrowbase) {(a)};
\node[panelmark] at (panelb.center |- panelrowbase) {(b)};
\node[panelmark] at (panelc.center |- panelrowbase) {(c)};
\end{tikzpicture}
